
\documentclass{article}
\usepackage{colortbl}
\usepackage{makecell}
\usepackage{multirow}
\usepackage{supertabular}

\begin{document}

\newcounter{utterance}

\centering \large Interaction Transcript for game `matchit\_ascii', experiment `similar\_grid\_1', episode 2 with gemma{-}4{-}26b{-}a4b{-}it.
\vspace{24pt}

{ \footnotesize  \setcounter{utterance}{1}
\setlength{\tabcolsep}{0pt}
\begin{supertabular}{c@{$\;$}|p{.15\linewidth}@{}p{.15\linewidth}p{.15\linewidth}p{.15\linewidth}p{.15\linewidth}p{.15\linewidth}}
    \# & $\;$A & \multicolumn{4}{c}{Game Master} & $\;\:$B\\
    \hline

    \theutterance \stepcounter{utterance}  
    & & \multicolumn{4}{p{0.6\linewidth}}{
        \cellcolor[rgb]{0.9,0.9,0.9}{
            \makecell[{{p{\linewidth}}}]{
                \texttt{\tiny{[P1$\langle$GM]}}
                \texttt{Estás participando en un juego colaborativo de adivinanzas.} \\
\texttt{El objetivo es averiguar si esta cuadrícula y otra imagen que solo yo puedo ver son iguales.} \\
\texttt{Por favor, describe primero tu cuadrícula. Luego, yo proporcionaré mi descripción y podremos hacernos preguntas sobre las cuadrículas para averiguar si son iguales. Ahora comienza tu breve descripción de la cuadrícula con "DESCRIPCIÓN:" seguida de la descripción. No agregues nada más. Esta es tu cuadrícula:} \\
\texttt{X X X X X} \\
\texttt{X ▢ ▢ ▢ ▢} \\
\texttt{X X X ▢ ▢} \\
\texttt{X ▢ ▢ ▢ ▢} \\
\texttt{X ▢ ▢ ▢ ▢} \\
            }
        }
    }
    & \\ \\

    \theutterance \stepcounter{utterance}  
    & \multicolumn{4}{p{0.6\linewidth}}{
        \cellcolor[rgb]{0.9,0.9,0.9}{
            \makecell[{{p{\linewidth}}}]{
                \texttt{\tiny{[P1$\rangle$GM]}}
                \texttt{DESCRIPCIÓN: Una cuadrícula de 5x5 donde la primera fila es toda X. La segunda fila tiene una X seguida de cuatro espacios vacíos. La tercera fila tiene dos X seguidas de dos espacios vacíos. La cuarta y quinta fila tienen una X seguida de cuatro espacios vacíos.} \\
            }
        }
    }
    & & \\ \\

    \theutterance \stepcounter{utterance}  
    & & & \multicolumn{2}{p{0.3\linewidth}}{
        \cellcolor[rgb]{0.9,0.9,0.9}{
            \makecell[{{p{\linewidth}}}]{
                \texttt{\tiny{[GM$|$GM]}}
                \texttt{continue} \\
            }
        }
    }
    & & \\ \\

    \theutterance \stepcounter{utterance}  
    & & \multicolumn{4}{p{0.6\linewidth}}{
        \cellcolor[rgb]{0.9,0.9,0.9}{
            \makecell[{{p{\linewidth}}}]{
                \texttt{\tiny{[GM$\rangle$P2]}}
                \texttt{Estás participando en un juego colaborativo de adivinanzas.} \\
\texttt{El objetivo es averiguar si esta cuadrícula y otra imagen que solo yo puedo ver son iguales.} \\
\texttt{Por favor, describe primero tu cuadrícula. Luego, yo proporcionaré mi descripción y podremos hacernos preguntas sobre las cuadrículas para averiguar si son iguales. Ahora comienza tu breve descripción de la cuadrícula con "DESCRIPCIÓN:" seguida de la descripción. No agregues nada más. Esta es tu cuadrícula:} \\
\texttt{X X X X X} \\
\texttt{▢ ▢ ▢ ▢ X} \\
\texttt{▢ ▢ X X X} \\
\texttt{▢ ▢ ▢ ▢ X} \\
\texttt{▢ ▢ ▢ ▢ X} \\
            }
        }
    }
    & \\ \\

    \theutterance \stepcounter{utterance}  
    & & & \multicolumn{4}{p{0.6\linewidth}}{
        \cellcolor[rgb]{0.9,0.9,0.9}{
            \makecell[{{p{\linewidth}}}]{
                \texttt{\tiny{[GM$\langle$P2]}}
                \texttt{DESCRIPCIÓN: Una cuadrícula de 5x5 donde la primera fila es toda "X". La segunda fila tiene cuatro "▢" y una "X" al final. La tercera fila tiene dos "▢" y tres "X". La cuarta y quinta fila tienen cuatro "▢" y una "X" al final.} \\
            }
        }
    }
     \\ \\

    \theutterance \stepcounter{utterance}  
    & & & \multicolumn{2}{p{0.3\linewidth}}{
        \cellcolor[rgb]{0.9,0.9,0.9}{
            \makecell[{{p{\linewidth}}}]{
                \texttt{\tiny{[GM$|$GM]}}
                \texttt{continue} \\
            }
        }
    }
    & & \\ \\

    \theutterance \stepcounter{utterance}  
    & & \multicolumn{4}{p{0.6\linewidth}}{
        \cellcolor[rgb]{0.9,0.9,0.9}{
            \makecell[{{p{\linewidth}}}]{
                \texttt{\tiny{[GM$\rangle$P2]}}
                \texttt{Esta es mi DESCRIPCIÓNDESCRIPCIÓN: Una cuadrícula de 5x5 donde la primera fila es toda X. La segunda fila tiene una X seguida de cuatro espacios vacíos. La tercera fila tiene dos X seguidas de dos espacios vacíos. La cuarta y quinta fila tienen una X seguida de cuatro espacios vacíos.} \\
\texttt{Ahora haz una pregunta para descubrir nuevos aspectos de mi cuadrícula que puedan ser diferentes a tu cuadrícula. Comienza con "PREGUNTA:" y no añadas nada más.} \\
            }
        }
    }
    & \\ \\

    \theutterance \stepcounter{utterance}  
    & & & \multicolumn{4}{p{0.6\linewidth}}{
        \cellcolor[rgb]{0.9,0.9,0.9}{
            \makecell[{{p{\linewidth}}}]{
                \texttt{\tiny{[GM$\langle$P2]}}
                \texttt{PREGUNTA: En tu segunda fila, ¿la "X" se encuentra en la primera posición (columna 1) o en la última posición (columna 5)?} \\
            }
        }
    }
     \\ \\

    \theutterance \stepcounter{utterance}  
    & & & \multicolumn{2}{p{0.3\linewidth}}{
        \cellcolor[rgb]{0.9,0.9,0.9}{
            \makecell[{{p{\linewidth}}}]{
                \texttt{\tiny{[GM$|$GM]}}
                \texttt{continue} \\
            }
        }
    }
    & & \\ \\

    \theutterance \stepcounter{utterance}  
    & & \multicolumn{4}{p{0.6\linewidth}}{
        \cellcolor[rgb]{0.9,0.9,0.9}{
            \makecell[{{p{\linewidth}}}]{
                \texttt{\tiny{[P1$\langle$GM]}}
                \texttt{Esta es mi DESCRIPCIÓNDESCRIPCIÓN: Una cuadrícula de 5x5 donde la primera fila es toda "X". La segunda fila tiene cuatro "▢" y una "X" al final. La tercera fila tiene dos "▢" y tres "X". La cuarta y quinta fila tienen cuatro "▢" y una "X" al final.} \\
\texttt{PREGUNTA: En tu segunda fila, ¿la "X" se encuentra en la primera posición (columna 1) o en la última posición (columna 5)?Comience su respuesta con "RESPUESTA:" y no añada nada más.} \\
            }
        }
    }
    & \\ \\

    \theutterance \stepcounter{utterance}  
    & \multicolumn{4}{p{0.6\linewidth}}{
        \cellcolor[rgb]{0.9,0.9,0.9}{
            \makecell[{{p{\linewidth}}}]{
                \texttt{\tiny{[P1$\rangle$GM]}}
                \texttt{RESPUESTA: En mi segunda fila, la "X" se encuentra en la primera posición (columna 1).} \\
            }
        }
    }
    & & \\ \\

    \theutterance \stepcounter{utterance}  
    & & & \multicolumn{2}{p{0.3\linewidth}}{
        \cellcolor[rgb]{0.9,0.9,0.9}{
            \makecell[{{p{\linewidth}}}]{
                \texttt{\tiny{[GM$|$GM]}}
                \texttt{continue} \\
            }
        }
    }
    & & \\ \\

    \theutterance \stepcounter{utterance}  
    & & \multicolumn{4}{p{0.6\linewidth}}{
        \cellcolor[rgb]{0.9,0.9,0.9}{
            \makecell[{{p{\linewidth}}}]{
                \texttt{\tiny{[P1$\langle$GM]}}
                \texttt{Ahora haz una pregunta para descubrir nuevos aspectos de mi cuadrícula que puedan ser diferentes a tu cuadrícula. Comienza con "PREGUNTA:" y no añadas nada más.} \\
            }
        }
    }
    & \\ \\

    \theutterance \stepcounter{utterance}  
    & \multicolumn{4}{p{0.6\linewidth}}{
        \cellcolor[rgb]{0.9,0.9,0.9}{
            \makecell[{{p{\linewidth}}}]{
                \texttt{\tiny{[P1$\rangle$GM]}}
                \texttt{PREGUNTA: En tu tercera fila, ¿la "X" se encuentra en las dos primeras posiciones (columnas 1 y 2) o en las tres últimas posiciones (columnas 3, 4 y 5)?} \\
            }
        }
    }
    & & \\ \\

    \theutterance \stepcounter{utterance}  
    & & & \multicolumn{2}{p{0.3\linewidth}}{
        \cellcolor[rgb]{0.9,0.9,0.9}{
            \makecell[{{p{\linewidth}}}]{
                \texttt{\tiny{[GM$|$GM]}}
                \texttt{continue} \\
            }
        }
    }
    & & \\ \\

    \theutterance \stepcounter{utterance}  
    & & \multicolumn{4}{p{0.6\linewidth}}{
        \cellcolor[rgb]{0.9,0.9,0.9}{
            \makecell[{{p{\linewidth}}}]{
                \texttt{\tiny{[GM$\rangle$P2]}}
                \texttt{RESPUESTA: En mi segunda fila, la "X" se encuentra en la primera posición (columna 1).} \\
\texttt{PREGUNTA: En tu tercera fila, ¿la "X" se encuentra en las dos primeras posiciones (columnas 1 y 2) o en las tres últimas posiciones (columnas 3, 4 y 5)?Comience su respuesta con "RESPUESTA:" y no añada nada más.} \\
            }
        }
    }
    & \\ \\

    \theutterance \stepcounter{utterance}  
    & & & \multicolumn{4}{p{0.6\linewidth}}{
        \cellcolor[rgb]{0.9,0.9,0.9}{
            \makecell[{{p{\linewidth}}}]{
                \texttt{\tiny{[GM$\langle$P2]}}
                \texttt{RESPUESTA: En mi tercera fila, las "X" se encuentran en las tres últimas posiciones (columnas 3, 4 y 5).} \\
            }
        }
    }
     \\ \\

    \theutterance \stepcounter{utterance}  
    & & & \multicolumn{2}{p{0.3\linewidth}}{
        \cellcolor[rgb]{0.9,0.9,0.9}{
            \makecell[{{p{\linewidth}}}]{
                \texttt{\tiny{[GM$|$GM]}}
                \texttt{continue} \\
            }
        }
    }
    & & \\ \\

    \theutterance \stepcounter{utterance}  
    & & \multicolumn{4}{p{0.6\linewidth}}{
        \cellcolor[rgb]{0.9,0.9,0.9}{
            \makecell[{{p{\linewidth}}}]{
                \texttt{\tiny{[GM$\rangle$P2]}}
                \texttt{Ahora haz una pregunta para descubrir nuevos aspectos de mi cuadrícula que puedan ser diferentes a tu cuadrícula. Comienza con "PREGUNTA:" y no añadas nada más.} \\
            }
        }
    }
    & \\ \\

    \theutterance \stepcounter{utterance}  
    & & & \multicolumn{4}{p{0.6\linewidth}}{
        \cellcolor[rgb]{0.9,0.9,0.9}{
            \makecell[{{p{\linewidth}}}]{
                \texttt{\tiny{[GM$\langle$P2]}}
                \texttt{PREGUNTA: En tu tercera fila, ¿las "X" están en las dos primeras posiciones (columnas 1 y 2) o en las tres últimas posiciones (columnas 3, 4 y 5)?} \\
            }
        }
    }
     \\ \\

    \theutterance \stepcounter{utterance}  
    & & & \multicolumn{2}{p{0.3\linewidth}}{
        \cellcolor[rgb]{0.9,0.9,0.9}{
            \makecell[{{p{\linewidth}}}]{
                \texttt{\tiny{[GM$|$GM]}}
                \texttt{continue} \\
            }
        }
    }
    & & \\ \\

    \theutterance \stepcounter{utterance}  
    & & \multicolumn{4}{p{0.6\linewidth}}{
        \cellcolor[rgb]{0.9,0.9,0.9}{
            \makecell[{{p{\linewidth}}}]{
                \texttt{\tiny{[P1$\langle$GM]}}
                \texttt{RESPUESTA: En mi tercera fila, las "X" se encuentran en las tres últimas posiciones (columnas 3, 4 y 5).} \\
\texttt{PREGUNTA: En tu tercera fila, ¿las "X" están en las dos primeras posiciones (columnas 1 y 2) o en las tres últimas posiciones (columnas 3, 4 y 5)?Comience su respuesta con "RESPUESTA:" y no añada nada más.} \\
            }
        }
    }
    & \\ \\

    \theutterance \stepcounter{utterance}  
    & \multicolumn{4}{p{0.6\linewidth}}{
        \cellcolor[rgb]{0.9,0.9,0.9}{
            \makecell[{{p{\linewidth}}}]{
                \texttt{\tiny{[P1$\rangle$GM]}}
                \texttt{RESPUESTA: En mi tercera fila, las "X" están en las dos primeras posiciones (columnas 1 y 2).} \\
            }
        }
    }
    & & \\ \\

    \theutterance \stepcounter{utterance}  
    & & & \multicolumn{2}{p{0.3\linewidth}}{
        \cellcolor[rgb]{0.9,0.9,0.9}{
            \makecell[{{p{\linewidth}}}]{
                \texttt{\tiny{[GM$|$GM]}}
                \texttt{continue} \\
            }
        }
    }
    & & \\ \\

    \theutterance \stepcounter{utterance}  
    & & \multicolumn{4}{p{0.6\linewidth}}{
        \cellcolor[rgb]{0.9,0.9,0.9}{
            \makecell[{{p{\linewidth}}}]{
                \texttt{\tiny{[P1$\langle$GM]}}
                \texttt{Ahora haz una pregunta para descubrir nuevos aspectos de mi cuadrícula que puedan ser diferentes a tu cuadrícula. Comienza con "PREGUNTA:" y no añadas nada más.} \\
            }
        }
    }
    & \\ \\

    \theutterance \stepcounter{utterance}  
    & \multicolumn{4}{p{0.6\linewidth}}{
        \cellcolor[rgb]{0.9,0.9,0.9}{
            \makecell[{{p{\linewidth}}}]{
                \texttt{\tiny{[P1$\rangle$GM]}}
                \texttt{PREGUNTA: En tu cuarta fila, ¿la "X" se encuentra en la primera posición (columna 1) o en la última posición (columna 5)?} \\
            }
        }
    }
    & & \\ \\

    \theutterance \stepcounter{utterance}  
    & & & \multicolumn{2}{p{0.3\linewidth}}{
        \cellcolor[rgb]{0.9,0.9,0.9}{
            \makecell[{{p{\linewidth}}}]{
                \texttt{\tiny{[GM$|$GM]}}
                \texttt{continue} \\
            }
        }
    }
    & & \\ \\

    \theutterance \stepcounter{utterance}  
    & & \multicolumn{4}{p{0.6\linewidth}}{
        \cellcolor[rgb]{0.9,0.9,0.9}{
            \makecell[{{p{\linewidth}}}]{
                \texttt{\tiny{[GM$\rangle$P2]}}
                \texttt{RESPUESTA: En mi tercera fila, las "X" están en las dos primeras posiciones (columnas 1 y 2).} \\
\texttt{PREGUNTA: En tu cuarta fila, ¿la "X" se encuentra en la primera posición (columna 1) o en la última posición (columna 5)?Comience su respuesta con "RESPUESTA:" y no añada nada más.} \\
            }
        }
    }
    & \\ \\

    \theutterance \stepcounter{utterance}  
    & & & \multicolumn{4}{p{0.6\linewidth}}{
        \cellcolor[rgb]{0.9,0.9,0.9}{
            \makecell[{{p{\linewidth}}}]{
                \texttt{\tiny{[GM$\langle$P2]}}
                \texttt{RESPUESTA: En mi cuarta fila, la "X" se encuentra en la última posición (columna 5).} \\
            }
        }
    }
     \\ \\

    \theutterance \stepcounter{utterance}  
    & & & \multicolumn{2}{p{0.3\linewidth}}{
        \cellcolor[rgb]{0.9,0.9,0.9}{
            \makecell[{{p{\linewidth}}}]{
                \texttt{\tiny{[GM$|$GM]}}
                \texttt{continue} \\
            }
        }
    }
    & & \\ \\

    \theutterance \stepcounter{utterance}  
    & & \multicolumn{4}{p{0.6\linewidth}}{
        \cellcolor[rgb]{0.9,0.9,0.9}{
            \makecell[{{p{\linewidth}}}]{
                \texttt{\tiny{[GM$\rangle$P2]}}
                \texttt{Ahora haz una pregunta para descubrir nuevos aspectos de mi cuadrícula que puedan ser diferentes a tu cuadrícula. Comienza con "PREGUNTA:" y no añadas nada más.} \\
            }
        }
    }
    & \\ \\

    \theutterance \stepcounter{utterance}  
    & & & \multicolumn{4}{p{0.6\linewidth}}{
        \cellcolor[rgb]{0.9,0.9,0.9}{
            \makecell[{{p{\linewidth}}}]{
                \texttt{\tiny{[GM$\langle$P2]}}
                \texttt{PREGUNTA: En tu quinta fila, ¿la "X" se encuentra en la primera posición (columna 1) o en la última posición (columna 5)?} \\
            }
        }
    }
     \\ \\

    \theutterance \stepcounter{utterance}  
    & & & \multicolumn{2}{p{0.3\linewidth}}{
        \cellcolor[rgb]{0.9,0.9,0.9}{
            \makecell[{{p{\linewidth}}}]{
                \texttt{\tiny{[GM$|$GM]}}
                \texttt{continue} \\
            }
        }
    }
    & & \\ \\

    \theutterance \stepcounter{utterance}  
    & & \multicolumn{4}{p{0.6\linewidth}}{
        \cellcolor[rgb]{0.9,0.9,0.9}{
            \makecell[{{p{\linewidth}}}]{
                \texttt{\tiny{[P1$\langle$GM]}}
                \texttt{RESPUESTA: En mi cuarta fila, la "X" se encuentra en la última posición (columna 5).} \\
\texttt{PREGUNTA: En tu quinta fila, ¿la "X" se encuentra en la primera posición (columna 1) o en la última posición (columna 5)?Comience su respuesta con "RESPUESTA:" y no añada nada más.} \\
            }
        }
    }
    & \\ \\

    \theutterance \stepcounter{utterance}  
    & \multicolumn{4}{p{0.6\linewidth}}{
        \cellcolor[rgb]{0.9,0.9,0.9}{
            \makecell[{{p{\linewidth}}}]{
                \texttt{\tiny{[P1$\rangle$GM]}}
                \texttt{RESPUESTA: En mi quinta fila, la "X" se encuentra en la primera posición (columna 1).} \\
            }
        }
    }
    & & \\ \\

    \theutterance \stepcounter{utterance}  
    & & & \multicolumn{2}{p{0.3\linewidth}}{
        \cellcolor[rgb]{0.9,0.9,0.9}{
            \makecell[{{p{\linewidth}}}]{
                \texttt{\tiny{[GM$|$GM]}}
                \texttt{continue} \\
            }
        }
    }
    & & \\ \\

    \theutterance \stepcounter{utterance}  
    & & \multicolumn{4}{p{0.6\linewidth}}{
        \cellcolor[rgb]{0.9,0.9,0.9}{
            \makecell[{{p{\linewidth}}}]{
                \texttt{\tiny{[P1$\langle$GM]}}
                \texttt{Ahora haz una pregunta para descubrir nuevos aspectos de mi cuadrícula que puedan ser diferentes a tu cuadrícula. Comienza con "PREGUNTA:" y no añadas nada más.} \\
            }
        }
    }
    & \\ \\

    \theutterance \stepcounter{utterance}  
    & \multicolumn{4}{p{0.6\linewidth}}{
        \cellcolor[rgb]{0.9,0.9,0.9}{
            \makecell[{{p{\linewidth}}}]{
                \texttt{\tiny{[P1$\rangle$GM]}}
                \texttt{PREGUNTA: En tu primera fila, ¿todas las posiciones son "X" o hay algún espacio vacío "▢"?} \\
            }
        }
    }
    & & \\ \\

    \theutterance \stepcounter{utterance}  
    & & & \multicolumn{2}{p{0.3\linewidth}}{
        \cellcolor[rgb]{0.9,0.9,0.9}{
            \makecell[{{p{\linewidth}}}]{
                \texttt{\tiny{[GM$|$GM]}}
                \texttt{continue} \\
            }
        }
    }
    & & \\ \\

    \theutterance \stepcounter{utterance}  
    & & \multicolumn{4}{p{0.6\linewidth}}{
        \cellcolor[rgb]{0.9,0.9,0.9}{
            \makecell[{{p{\linewidth}}}]{
                \texttt{\tiny{[GM$\rangle$P2]}}
                \texttt{RESPUESTA: En mi quinta fila, la "X" se encuentra en la primera posición (columna 1).} \\
\texttt{PREGUNTA: En tu primera fila, ¿todas las posiciones son "X" o hay algún espacio vacío "▢"?Comience su respuesta con "RESPUESTA:" y no añada nada más.} \\
            }
        }
    }
    & \\ \\

    \theutterance \stepcounter{utterance}  
    & & & \multicolumn{4}{p{0.6\linewidth}}{
        \cellcolor[rgb]{0.9,0.9,0.9}{
            \makecell[{{p{\linewidth}}}]{
                \texttt{\tiny{[GM$\langle$P2]}}
                \texttt{RESPUESTA: En mi primera fila, todas las posiciones son "X".} \\
            }
        }
    }
     \\ \\

    \theutterance \stepcounter{utterance}  
    & & & \multicolumn{2}{p{0.3\linewidth}}{
        \cellcolor[rgb]{0.9,0.9,0.9}{
            \makecell[{{p{\linewidth}}}]{
                \texttt{\tiny{[GM$|$GM]}}
                \texttt{continue} \\
            }
        }
    }
    & & \\ \\

    \theutterance \stepcounter{utterance}  
    & & \multicolumn{4}{p{0.6\linewidth}}{
        \cellcolor[rgb]{0.9,0.9,0.9}{
            \makecell[{{p{\linewidth}}}]{
                \texttt{\tiny{[GM$\rangle$P2]}}
                \texttt{Ahora toma una decisión. ¿Qué piensas: tu cuadrícula y la otra cuadrícula descrita son la misma? Escribe "DECISIÓN: misma cuadrícula" si crees que son la misma cuadrícula o "DECISIÓN: cuadrículas diferentes" si crees que son cuadrículas diferentes. No añadas nada más.} \\
            }
        }
    }
    & \\ \\

    \theutterance \stepcounter{utterance}  
    & & & \multicolumn{4}{p{0.6\linewidth}}{
        \cellcolor[rgb]{0.9,0.9,0.9}{
            \makecell[{{p{\linewidth}}}]{
                \texttt{\tiny{[GM$\langle$P2]}}
                \texttt{DECISIÓN: cuadrículas diferentes} \\
            }
        }
    }
     \\ \\

    \theutterance \stepcounter{utterance}  
    & & & \multicolumn{2}{p{0.3\linewidth}}{
        \cellcolor[rgb]{0.9,0.9,0.9}{
            \makecell[{{p{\linewidth}}}]{
                \texttt{\tiny{[GM$|$GM]}}
                \texttt{continue} \\
            }
        }
    }
    & & \\ \\

    \theutterance \stepcounter{utterance}  
    & & & \multicolumn{2}{p{0.3\linewidth}}{
        \cellcolor[rgb]{0.9,0.9,0.9}{
            \makecell[{{p{\linewidth}}}]{
                \texttt{\tiny{[GM$|$GM]}}
                \texttt{success} \\
            }
        }
    }
    & & \\ \\

    \theutterance \stepcounter{utterance}  
    & & \multicolumn{4}{p{0.6\linewidth}}{
        \cellcolor[rgb]{0.9,0.9,0.9}{
            \makecell[{{p{\linewidth}}}]{
                \texttt{\tiny{[P1$\langle$GM]}}
                \texttt{RESPUESTA: En mi primera fila, todas las posiciones son "X".} \\
\texttt{Ahora toma una decisión. ¿Qué piensas: tu cuadrícula y la otra cuadrícula descrita son la misma? Escribe "DECISIÓN: misma cuadrícula" si crees que son la misma cuadrícula o "DECISIÓN: cuadrículas diferentes" si crees que son cuadrículas diferentes. No añadas nada más.} \\
            }
        }
    }
    & \\ \\

    \theutterance \stepcounter{utterance}  
    & \multicolumn{4}{p{0.6\linewidth}}{
        \cellcolor[rgb]{0.9,0.9,0.9}{
            \makecell[{{p{\linewidth}}}]{
                \texttt{\tiny{[P1$\rangle$GM]}}
                \texttt{DECISIÓN: cuadrículas diferentes} \\
            }
        }
    }
    & & \\ \\

    \theutterance \stepcounter{utterance}  
    & & & \multicolumn{2}{p{0.3\linewidth}}{
        \cellcolor[rgb]{0.9,0.9,0.9}{
            \makecell[{{p{\linewidth}}}]{
                \texttt{\tiny{[GM$|$GM]}}
                \texttt{continue} \\
            }
        }
    }
    & & \\ \\

    \theutterance \stepcounter{utterance}  
    & & & \multicolumn{2}{p{0.3\linewidth}}{
        \cellcolor[rgb]{0.9,0.9,0.9}{
            \makecell[{{p{\linewidth}}}]{
                \texttt{\tiny{[GM$|$GM]}}
                \texttt{success} \\
            }
        }
    }
    & & \\ \\

\end{supertabular}
}

\end{document}
